\documentclass[aspectratio=169]{beamer}

\usetheme{Madrid}
\usecolortheme{default}

\usepackage{amsmath, amssymb}
\usepackage{booktabs}
\usepackage{array}
\usepackage{xcolor}

\title[Barcoding Detection]{Barcoding Detection and Parameter Sensitivity}
\subtitle{OCT Profile Extraction Proof-of-Concept}
\author{Jonathan Ma}
\date{July 31, 2026}

\begin{document}

\begin{frame}
  \titlepage
\end{frame}

\begin{frame}{Roadmap}
  \tableofcontents
\end{frame}

\section{Interactive Profile Overlay and Manual Classification}

\begin{frame}{New Tool: Interactive Profile Overlay}

\begin{block}{Purpose}
The annotation interface aligns the processed sub-BM scan with its extracted
intensity profile on a shared horizontal axis.
\end{block}

\begin{columns}[T]

\column{0.50\textwidth}

\textbf{Workflow}

\begin{enumerate}
    \item Select an annotation class:
    \begin{itemize}
        \item Normal
        \item Early Atrophy (EA)
        \item Barcoding
    \end{itemize}
    \item Drag horizontally across either:
    \begin{itemize}
        \item the processed OCT scan, or
        \item the aligned intensity profile.
    \end{itemize}
    \item Use \texttt{Undo last} or \texttt{Clear all} to revise intervals.
    \item Save the annotations when enabled.
\end{enumerate}

\column{0.46\textwidth}

\textbf{Saved information}

\begin{itemize}
    \item Source E2E file
    \item Central B-scan index
    \item Profile depth and step size
    \item Profile margin and aggregation method
    \item Normalization configuration
    \item Interval start and end positions
    \item Class label and interval width
\end{itemize}

\end{columns}

\end{frame}


\begin{frame}{Manual Classification Results}

\begin{block}{Initial proof-of-concept annotation}
Three contiguous horizontal regions were manually classified on the central
B-scan.
\end{block}

\begin{table}
\centering
\small
\begin{tabular}{lccc}
\toprule
\textbf{Class} &
\textbf{Start $x$} &
\textbf{End $x$} &
\textbf{Width (px)} \\
\midrule
Normal &
2.0 &
131.5 &
129.5 \\
Early Atrophy (EA) &
137.1 &
185.0 &
47.9 \\
Barcoding &
191.9 &
263.5 &
71.6 \\
\bottomrule
\end{tabular}
\end{table}

\vspace{0.3cm}

\begin{columns}[T]

\column{0.50\textwidth}

\textbf{Interval counts}

\begin{itemize}
    \item Normal: 1
    \item Early Atrophy: 1
    \item Barcoding: 1
\end{itemize}

\column{0.46\textwidth}

\textbf{Total annotated width}

\begin{itemize}
    \item Normal: 129.5 px
    \item Early Atrophy: 47.9 px
    \item Barcoding: 71.6 px
\end{itemize}

\end{columns}

\end{frame}

\section{Profile Extraction Tests}

\begin{frame}{Profile Extraction Expansion}

We want to evaluate whether averaging neighboring rows produces a more robust intensity
profile than extracting a single horizontal line.


\textbf{Configurations evaluated}

\begin{table}
\centering
\small
\begin{tabular}{ll}
\toprule
Method & Description\\
\midrule
Single Line & Baseline profile from one row\\
Mean $\pm2$ px & Average across 5 rows\\
Mean $\pm4$ px & Average across 9 rows\\
Median $\pm2$ px & Median across 5 rows\\
Median $\pm4$ px & Median across 9 rows\\
Median $\pm6$ px & Median across 13 rows\\
\bottomrule
\end{tabular}
\end{table}


\textbf{Summary statistics}

\begin{itemize}
\item Gradient energy — amount of rapid intensity change, profile roughness.
\item Correlation — similarity to the original single-line extraction (1.0 = identical).
\item (MAD) — average point-by-point intensity difference from the baseline.
\item (RMSD) — overall magnitude of deviation, emphasizing larger differences.
\end{itemize}

\end{frame}

\begin{frame}{Profile Extraction Results}

\small

\begin{table}
\centering
\begin{tabular}{lcccc}
\toprule
Method &
Gradient &
Correlation &
MAD &
RMSD\\
\midrule
Single Line & 270.8 & 1.000 & 0.00 & 0.00\\
Mean $\pm2$ & 108.0 & 0.933 & 7.45 & 9.27\\
Mean $\pm4$ & 81.4 & 0.908 & 8.60 & 10.76\\
Median $\pm2$ & 146.3 & 0.935 & 6.15 & 9.11\\
Median $\pm4$ & 107.2 & 0.907 & 7.82 & 10.76\\
Median $\pm6$ & 83.5 & 0.901 & 8.35 & 11.11\\
\bottomrule
\end{tabular}
\end{table}

\textbf{Observations}

\begin{itemize}
\item Increasing the averaging band substantially reduced roughness.
\item Mean and median aggregation produced similar overall profiles.
\item Median $\pm2$ produced the smallest deviation from the baseline profile.
\item Larger averaging windows generated smoother profiles but differed more.
\end{itemize}

\end{frame}

\begin{frame}{Interpretation}

\begin{block}{Clinical interpretation}

\begin{itemize}
\item Averaging neighboring rows suppresses local speckle noise while preserving the overall profile shape.
\item Small averaging windows ($\pm2$ pixels) appear to provide the best balance between noise reduction and preservation of local features.
\item Very large averaging windows may smooth clinically relevant local intensity changes associated with early atrophy or barcoding.
\end{itemize}

\end{block}

\vspace{0.4cm}

\begin{alertblock}{Current recommendation}

For subsequent experiments, a profile band of approximately
$\pm2$ pixels is recommended as the baseline extraction strategy.
This configuration substantially reduces profile variability while
remaining highly correlated ($r\approx0.93$) with the original
single-line extraction.

\end{alertblock}

\end{frame}

\section{Normalization Scopes}

\begin{frame}{Normalization Scope Comparison}

\textbf{Objective}

Evaluate whether normalization should be performed using the entire
sub-Bruch's membrane (BM) region or only the rows used for profile extraction.

\vspace{0.3cm}

\textbf{Normalization strategies}

\begin{table}
\centering
\small
\begin{tabular}{ll}
\toprule
Method & Description\\
\midrule
Whole ROI & Percentiles computed using all pixels below BM\\
Local Band & Percentiles computed only within the extraction band\\
\bottomrule
\end{tabular}
\end{table}

\vspace{0.3cm}

Both strategies were evaluated using:

\begin{itemize}
\item Robust percentile normalization (1--99\%) \& Min--max normalization (0--100\%)
\end{itemize}

\vspace{0.2cm}

\textbf{Summary statistics}

\begin{itemize}
\item Correlation, MAD, RMSD, Gradient Energy, and dynamic range, the difference between the darkest and brightest profile values after normalization.
\end{itemize}

\end{frame}

\begin{frame}{Normalization Scope Results}

\small

\begin{table}
\centering
\begin{tabular}{lcccc}
\toprule
Method &
Correlation &
MAD &
RMSD &
Gradient\\
\midrule
Whole ROI (1--99) & 1.0000 & 0.00 & 0.00 & 180.1\\
Local Band (1--99) & 0.9971 & 10.24 & 12.49 & 328.3\\
Whole ROI (0--100) & 0.99996 & 39.80 & 41.26 & 84.0\\
Local Band (0--100) & 0.99996 & 10.67 & 11.05 & 151.4\\
\bottomrule
\end{tabular}
\end{table}

\vspace{0.3cm}

\textbf{Normalization intensity ranges}

\begin{table}
\centering
\small
\begin{tabular}{lcc}
\toprule
Method & Lower & Upper\\
\midrule
Whole ROI (1--99) & 0.0 & 171.0\\
Local Band (1--99) & 25.1 & 149.0\\
Whole ROI (0--100) & 0.0 & 251.0\\
Local Band (0--100) & 0.0 & 187.0\\
\bottomrule
\end{tabular}
\end{table}

\end{frame}

\begin{frame}{Interpretation}

\begin{block}{Clinical interpretation}

\begin{itemize}
\item Whole-ROI normalization provides a consistent intensity reference using the entire processed scan.
\item Local-band normalization adapts the intensity scaling to the immediate neighborhood surrounding the extracted profile.
\item Restricting normalization to the local band increases local contrast, producing larger intensity gradients that may improve visibility of subtle features.
\item Min--max normalization (0--100\%) uses the full intensity range and is more influenced by extreme pixel values than percentile-based normalization.
\end{itemize}

\end{block}

\vspace{0.4cm}

\begin{alertblock}{Current recommendation}

Percentile normalization  using the entire processed ROI is recommended because it provides a stable global reference while remaining robust to extreme intensity values.

\end{alertblock}

\end{frame}

\section{One-Factor-at-a-Time Analysis}

\begin{frame}{One-Factor-at-a-Time Experimental Setup}

\textbf{Objective}

Measure how strongly each processing choice changes the extracted intensity
profile when all other settings are held fixed.

\vspace{0.25cm}

\textbf{Baseline configuration}

\begin{table}
\centering
\small
\begin{tabular}{ll}
\toprule
Parameter & Baseline setting\\
\midrule
Profile margin & $\pm4$ pixels\\
Aggregation & Mean\\
Normalization scope & Whole ROI\\
Percentile bounds & 1--99\%\\
Sampling step size & 1 pixel\\
\bottomrule
\end{tabular}
\end{table}

\vspace{0.2cm}

\textbf{Factors varied independently}

\begin{itemize}
    \item Profile margin: $0$, $\pm2$, $\pm4$, $\pm6$, and $\pm8$ pixels
    \item Percentile bounds: 0--100, 0.5--99.5, 1--99, 2--98, and 5--95
    \item Sampling step size: 1, 2, 4, and 8 pixels
\end{itemize}

\end{frame}

\begin{frame}{How OFAT Sensitivity Was Measured}

\small

Each alternative profile was compared with the baseline profile at the same
horizontal image positions.

\begin{itemize}
    \item \textbf{Sufficient Statistics}: Correlation, MAD, RMSD, Max Absolute Difference, Grad Energy

    \item \textbf{Curvature energy}: the amount of rapid change in profile
    slope.
    \begin{itemize}
        \item Helps measure sharp peaks and narrow local features.
    \end{itemize}
\end{itemize}

\end{frame}

\begin{frame}{OFAT Results: Profile Margin and Aggregation}

\small

\begin{table}
\centering
\begin{tabular}{lcccc}
\toprule
Setting &
Correlation &
MAD &
RMSD &
Roughness\\
\midrule
Margin $\pm0$ px & 0.907 & 12.80 & 16.00 & 207.2\\
Margin $\pm2$ px & 0.986 & 4.68 & 5.88 & 103.1\\
Margin $\pm4$ px & 1.000 & 0.00 & 0.00 & 83.8\\
Margin $\pm6$ px & 0.992 & 3.44 & 4.39 & 68.2\\
Margin $\pm8$ px & 0.983 & 4.77 & 6.26 & 62.3\\
\midrule
Mean aggregation & 1.000 & 0.00 & 0.00 & 83.8\\
Median aggregation & 0.989 & 4.01 & 5.14 & 99.4\\
\bottomrule
\end{tabular}
\end{table}

\vspace{0.25cm}

\textbf{Interpretation}

\begin{itemize}
    \item The single-line profile was substantially rougher and differed most
    from the baseline.
    \item Increasing the margin progressively smoothed the profile.
    \item Margins of $\pm2$ to $\pm6$ pixels remained highly similar to the
    baseline.
    \item Mean and median aggregation produced similar profile shapes, although
    the median retained slightly more local roughness.
\end{itemize}

\end{frame}

\begin{frame}{OFAT Results: Sampling Step Size}

\small

\begin{table}
\centering
\begin{tabular}{lcccc}
\toprule
Step size &
Correlation &
MAD &
RMSD &
Roughness\\
\midrule
1 px & 1.000 & 0.00 & 0.00 & 83.8\\
2 px & 0.981 & 3.55 & 6.71 & 66.3\\
4 px & 0.954 & 6.95 & 10.36 & 34.7\\
8 px & 0.894 & 11.09 & 15.95 & 13.5\\
\bottomrule
\end{tabular}
\end{table}

\vspace{0.3cm}

\begin{block}{Interpretation}

\begin{itemize}
    \item Larger step sizes reduce the number of sampled horizontal positions.
    \item This produces progressively smoother profiles because narrow local
    changes may be skipped.
    \item A 2-pixel step remained relatively close to the baseline.
    \item Step sizes of 4 or 8 pixels increasingly altered the profile shape
    and may miss narrow barcode-like features.
\end{itemize}

\end{block}

\end{frame}

\begin{frame}{Which Parameters Had the Greatest Influence?}

\small

\begin{table}
\centering
\begin{tabular}{lccc}
\toprule
Factor &
Mean MAD &
Maximum RMSD &
Minimum correlation\\
\midrule
Percentile bounds & 15.80 & 41.26 & 0.998\\
Sampling step size & 5.40 & 15.95 & 0.894\\
Profile margin & 5.14 & 16.00 & 0.907\\
Normalization scope & 5.12 & 12.49 & 0.997\\
Aggregation method & 2.01 & 5.14 & 0.989\\
\bottomrule
\end{tabular}
\end{table}

\vspace{0.3cm}

\textbf{Clinical interpretation}

\begin{itemize}
    \item \textbf{Percentile bounds} changed the absolute gray-value scale most
    strongly, although the spatial profile shape was largely preserved.
    \item \textbf{Sampling step size} and \textbf{profile margin} had the
    greatest effect on profile shape and local detail.
    \item \textbf{Aggregation method} had the smallest overall influence.
\end{itemize}

\end{frame}

\begin{frame}{OFAT Conclusions}

\begin{alertblock}{Current working configuration}

\begin{itemize}
    \item Profile margin: approximately $\pm4$ pixels
    \item Aggregation: mean
    \item Normalization: whole-ROI, 1--99\%
    \item Sampling step size: 1 pixel
\end{itemize}

\end{alertblock}

\vspace{0.3cm}

\begin{itemize}
    \item Small symmetric averaging bands reduce speckle while preserving the
    underlying spatial profile.
    \item Whole-ROI percentile normalization provides a reproducible intensity
    reference.
    \item Native 1-pixel sampling is preferred for preserving narrow barcode
    structures.
    \item These findings are based on the current proof-of-concept scan and
    require confirmation across additional OCT volumes.
\end{itemize}

\end{frame}

\section{Factorial Analysis}

\begin{frame}{Factorial Experimental Setup}

\textbf{Objective}

Evaluate combinations of preprocessing and profile-extraction settings to
determine whether parameter interactions produce substantially different
intensity profiles.

\vspace{0.25cm}

\textbf{Factors included}

\begin{table}
\centering
\small
\begin{tabular}{ll}
\toprule
Parameter & Levels evaluated\\
\midrule
Profile margin & $\pm2$, $\pm4$, and $\pm6$ pixels\\
Aggregation & Mean and median\\
Normalization scope & Whole ROI and local band\\
Percentile bounds & 0--100, 1--99, and 2--98\%\\
Sampling step size & 1 and 2 pixels\\
\bottomrule
\end{tabular}
\end{table}

\vspace{0.25cm}

\[
3 \times 2 \times 2 \times 3 \times 2
=
72 \text{ parameter combinations}
\]

Each resulting profile was compared with the working baseline:

\begin{center}
$\pm4$ pixels, mean aggregation, whole-ROI 1--99\% normalization,
and 1-pixel sampling.
\end{center}

\end{frame}


\begin{frame}{Most Stable Factorial Configurations}

\scriptsize

\begin{table}
\centering
\begin{tabular}{cccccccc}
\toprule
Margin &
Aggregate &
Scope &
Bounds &
Step &
Corr. &
MAD &
RMSD\\
\midrule
$\pm4$ & Mean & Whole ROI & 1--99 & 1 & 1.000 & 0.00 & 0.00\\
$\pm6$ & Mean & Whole ROI & 1--99 & 1 & 0.992 & 3.44 & 4.39\\
$\pm4$ & Median & Whole ROI & 1--99 & 1 & 0.989 & 4.01 & 5.14\\
$\pm6$ & Median & Whole ROI & 1--99 & 1 & 0.987 & 4.47 & 5.59\\
$\pm2$ & Mean & Whole ROI & 1--99 & 1 & 0.986 & 4.68 & 5.88\\
$\pm4$ & Mean & Whole ROI & 2--98 & 1 & 1.000 & 6.09 & 6.30\\
$\pm4$ & Mean & Whole ROI & 1--99 & 2 & 0.981 & 3.55 & 6.71\\
\bottomrule
\end{tabular}
\end{table}

\vspace{0.25cm}

\textbf{Observed pattern}

\begin{itemize}
    \item The most stable configurations consistently used whole-ROI
    normalization.
    \item Percentile bounds of 1--99\% appeared most frequently among the
    highest-ranked settings.
    \item Margins of $\pm4$ or $\pm6$ pixels were more stable than the smallest
    tested margin.
    \item One-pixel sampling generally produced the closest agreement with the
    baseline.
\end{itemize}

\end{frame}

\begin{frame}{Configurations Producing the Largest Changes}

\scriptsize

\begin{table}
\centering
\begin{tabular}{cccccccc}
\toprule
Margin &
Aggregate &
Scope &
Bounds &
Step &
Corr. &
MAD &
RMSD\\
\midrule
$\pm6$ & Mean & Whole ROI & 0--100 & 1 & 0.992 & 40.06 & 41.79\\
$\pm6$ & Mean & Whole ROI & 0--100 & 2 & 0.976 & 39.66 & 41.71\\
$\pm2$ & Mean & Whole ROI & 0--100 & 2 & 0.970 & 39.77 & 41.70\\
$\pm6$ & Median & Whole ROI & 0--100 & 1 & 0.987 & 39.87 & 41.65\\
$\pm6$ & Median & Whole ROI & 0--100 & 2 & 0.973 & 39.60 & 41.64\\
$\pm2$ & Median & Whole ROI & 0--100 & 2 & 0.961 & 39.45 & 41.48\\
\bottomrule
\end{tabular}
\end{table}

\vspace{0.25cm}

\textbf{Interpretation}

\begin{itemize}
    \item Min--max normalization produced the largest gray-value shifts.
    \item Correlations remained relatively high, indicating that the overall
    spatial pattern was preserved despite substantial scale changes.
    \item This means two profiles can have similar shape while having very
    different absolute intensity values.
    \item Local-band 2--98\% normalization also produced larger deviations,
    although less extreme than whole-ROI min--max normalization.
\end{itemize}

\end{frame}

\begin{frame}{Balancing Stability and Signal Preservation}

\small

Stability alone is not sufficient: excessive smoothing may reduce noise while
also suppressing clinically relevant barcode-like intensity changes.

\begin{itemize}
    \item \textbf{Gradient energy ratio}
    \begin{itemize}
        \item A value near 1 indicates similar preservation of local intensity
        changes.
        \item Values below 1 indicate stronger smoothing.
        \item Values above 1 indicate greater local variation.
    \end{itemize}

    \item \textbf{Stability--signal score}
    \begin{itemize}
        \item Lower scores indicate better agreement with the baseline while
        retaining a similar amount of local structure.
        \item This is an exploratory ranking measure rather than a clinical
        diagnostic score.
    \end{itemize}
\end{itemize}

\begin{table}
\centering
\scriptsize
\begin{tabular}{lccc}
\toprule
Configuration & RMSD & Gradient ratio & Score\\
\midrule
$\pm4$, mean, whole ROI, 1--99, step 1 & 0.00 & 1.000 & -2.577\\
$\pm6$, median, whole ROI, 1--99, step 1 & 5.59 & 0.974 & -2.050\\
$\pm4$, mean, whole ROI, 2--98, step 1 & 6.30 & 1.094 & -1.854\\
$\pm6$, mean, whole ROI, 1--99, step 1 & 4.39 & 0.814 & -1.829\\
$\pm4$, median, whole ROI, 1--99, step 1 & 5.14 & 1.186 & -1.766\\
\bottomrule
\end{tabular}
\end{table}

\end{frame}

\begin{frame}{Factorial Analysis Conclusions}

\begin{block}{Main findings}

\begin{itemize}
    \item Stable parameter combinations clustered around whole-ROI percentile
    normalization, small-to-moderate averaging bands, and native 1-pixel
    sampling.
    \item Percentile bounds strongly affected absolute intensity scale.
    \item Profile margin and sampling step size more directly affected local
    profile detail.
    \item Mean and median aggregation were both viable, with comparatively
    smaller effects than normalization and sampling choices.
\end{itemize}

\end{block}

\vspace{0.25cm}

\begin{alertblock}{Working configuration}
Use 0.01-0.99 normalization about the whole ROI, using mean $\pm4$ pixels

\end{alertblock}

\vspace{0.2cm}

\centering
\textit{These results describe one proof-of-concept scan and will be
validated across the remaining OCT volumes before fixing final defaults.}

\end{frame}



\section{Conclusions}

\begin{frame}{Conclusions}

\begin{block}{Main findings}

\begin{itemize}
    \item Small-to-moderate depth averaging reduced profile roughness while preserving overall signal shape.
    \item Whole-ROI percentile normalization provided the most stable intensity reference.
    \item Sampling step size strongly affected preservation of narrow local features.
    \item Aggregation method had a smaller influence than normalization, margin, or step size.
    \item Stable configurations clustered around:
    \begin{itemize}
        \item $\pm4$ to $\pm6$ pixel margins,
        \item whole-ROI 1--99\% normalization,
        \item and 1-pixel sampling.
    \end{itemize}
\end{itemize}

\end{block}

\end{frame}

\begin{frame}{Next Steps}

\begin{enumerate}
    \item \textbf{Apply the pipeline to all available E2E volumes}
    \begin{itemize}
        \item Test the central scan of each volume
        \item Test across all scans of a single volume
    \end{itemize}

    \item \textbf{Develop numerical detection features}
    \begin{itemize}
        \item Local intensity level
        \item Gradient and curvature activity
        \item Peak density and spatial frequency
        \item Vertical persistence across neighboring depths
        \item Look into Neural Super Resolution (a resolution enhancing technique like DLSS)
    \end{itemize}

    \item \textbf{Evaluate automated spatial classification}
    \begin{itemize}
        \item Threshold-based signal processing
        \item Change-point methods
        \item Hidden Markov models
    \end{itemize}

\end{enumerate}

\end{frame}

\end{document}